\documentclass[12pt,letterpaper]{article}
\usepackage{graphicx,textcomp}
\usepackage{natbib}
\usepackage{setspace}
\usepackage{fullpage}
\usepackage{color}
\usepackage[reqno]{amsmath}
\usepackage{amsthm}
\usepackage{fancyvrb}
\usepackage{amssymb,enumerate}
\usepackage[all]{xy}
\usepackage{endnotes}
\usepackage{lscape}
\newtheorem{com}{Comment}
\usepackage{float}
\usepackage{hyperref}
\newtheorem{lem} {Lemma}
\newtheorem{prop}{Proposition}
\newtheorem{thm}{Theorem}
\newtheorem{defn}{Definition}
\newtheorem{cor}{Corollary}
\newtheorem{obs}{Observation}
\usepackage[compact]{titlesec}
\usepackage{dcolumn}
\usepackage{tikz}
\usetikzlibrary{arrows}
\usepackage{multirow}
\usepackage{xcolor}
\newcolumntype{.}{D{.}{.}{-1}}
\newcolumntype{d}[1]{D{.}{.}{#1}}
\definecolor{light-gray}{gray}{0.65}
\usepackage{url}
\usepackage{listings}
\usepackage{color}
\usepackage{booktabs}


\definecolor{codegreen}{rgb}{0,0.6,0}
\definecolor{codegray}{rgb}{0.5,0.5,0.5}
\definecolor{codepurple}{rgb}{0.58,0,0.82}
\definecolor{backcolour}{rgb}{0.95,0.95,0.92}

\lstdefinestyle{mystyle}{
	backgroundcolor=\color{backcolour},   
	commentstyle=\color{codegreen},
	keywordstyle=\color{magenta},
	numberstyle=\tiny\color{codegray},
	stringstyle=\color{codepurple},
	basicstyle=\footnotesize,
	breakatwhitespace=false,         
	breaklines=true,                 
	captionpos=b,                    
	keepspaces=true,                 
	numbers=left,                    
	numbersep=5pt,                  
	showspaces=false,                
	showstringspaces=false,
	showtabs=false,                  
	tabsize=2
}
\lstset{style=mystyle}
\newcommand{\Sref}[1]{Section~\ref{#1}}
\newtheorem{hyp}{Hypothesis}

\title{Problem Set 2}
\date{Due: February 18, 2026}
\author{Applied Stats II}


\begin{document}
	\maketitle
	\section*{Instructions}
	\begin{itemize}
		\item Please show your work! You may lose points by simply writing in the answer. If the problem requires you to execute commands in \texttt{R}, please include the code you used to get your answers. Please also include the \texttt{.R} file that contains your code. If you are not sure if work needs to be shown for a particular problem, please ask.
		\item Your homework should be submitted electronically on GitHub in \texttt{.pdf} form.
		\item This problem set is due before 23:59 on Wednesday February 18, 2026. No late assignments will be accepted.

	\end{itemize}

	\vspace{.25cm}

\noindent We're interested in what types of international environmental agreements or policies people support (\href{https://www.pnas.org/content/110/34/13763}{Bechtel and Scheve 2013)}. So, we asked 8,500 individuals whether they support a given policy, and for each participant, we vary the (1) number of countries that participate in the international agreement and (2) sanctions for not following the agreement. \\

\noindent Load in the data labeled \texttt{climateSupport.RData} on GitHub, which contains an observational study of 8,500 observations.

\begin{itemize}
	\item
	Response variable: 
	\begin{itemize}
		\item \texttt{choice}: 1 if the individual agreed with the policy; 0 if the individual did not support the policy
	\end{itemize}
	\item
	Explanatory variables: 
	\begin{itemize}
		\item
		\texttt{countries}: Number of participating countries [20 of 192; 80 of 192; 160 of 192]
		\item
		\texttt{sanctions}: Sanctions for missing emission reduction targets [None, 5\%, 15\%, and 20\% of the monthly household costs given 2\% GDP growth]
		
	\end{itemize}
	
\end{itemize}

\newpage
\noindent Please answer the following questions:

\begin{enumerate}
	\item
	Remember, we are interested in predicting the likelihood of an individual supporting a policy based on the number of countries participating and the possible sanctions for non-compliance.
	\begin{enumerate}
		\item [] Fit an additive model. Provide the summary output, the global null hypothesis, and $p$-value. Please describe the results and provide a conclusion.
		\lstinputlisting[language=R, firstline=1,lastline=67]{PS02_RS.R} 
			\vspace{10cm}
		\begin{table}[htbp]
			\centering
			\caption{Logistic Regression Results: Climate Support}
			\begin{tabular}{lrrrr}
				\toprule
				& Estimate & Std. Error & z value & Pr($>|z|$) \\
				\midrule
				(Intercept)   & -0.0057 & 0.0220 & -0.258 & 0.7965 \\
				countries.L   & 0.4585  & 0.0381 & 12.033 & $< 2\times10^{-16}$*** \\
				countries.Q   & -0.0100 & 0.0381 & -0.261 & 0.7937 \\
				sanctions.L   & -0.2763 & 0.0439 & -6.291 & $3.15\times10^{-10}$*** \\
				sanctions.Q   & -0.1811 & 0.0440 & -4.119 & $3.80\times10^{-5}$*** \\
				sanctions.C   & 0.1502  & 0.0440 & 3.414  & 0.000639*** \\
				\midrule
				\multicolumn{5}{l}{\footnotesize Signif. codes: *** $p<0.001$, ** $p<0.01$, * $p<0.05$} \\
				\bottomrule
			\end{tabular}
			
			\vspace{0.3cm}
			
	\begin{center}
		\begin{tabular}{lcccccc}
			\hline
			Model & Resid.\ Df & Resid.\ Dev & Df & Deviance & Pr($>$Chi) \\
			\hline
			1: choice $\sim$ 1 & 8499 & 11783 &  &  &  \\
			2: choice $\sim$ countries + sanctions & 8494 & 11568 & 5 & 215.15 & $< 2.2\times10^{-16}$*** \\
			\hline
		\end{tabular}
	\end{center}
	\vspace{0.3cm}
			
			\begin{tabular}{ll}
				Model & Logistic regression (binomial link) \\
				Observations & 8500 \\
				Null deviance & 11783 (df = 8499) \\
				Residual deviance & 11568 (df = 8494) \\
				AIC & 11580 \\
				Fisher scoring iterations & 4 \\
			\end{tabular}
			
		\end{table}
		
An additive logistic regression model was fitted to predict whether an individual supports an international environmental policy using the number of participating countries and the level of sanctions for non-compliance as explanatory variables.
	
	The global null hypothesis for this model is:
	
	\[
	H_0: \beta_{\text{countries}} = \beta_{\text{sanctions}} = 0
	\]
	
	
	which states that neither the number of participating countries nor the level of sanctions has an effect on the likelihood of supporting the policy. The alternative hypothesis is that at least one coefficient differs from zero.
		\vspace{0.3cm}
	The likelihood ratio test comparing the additive model to the intercept-only model produces a very small p-value (p < 0.05), indicating that the model with explanatory variables provides a significantly better fit than the null model. Therefore, we reject the global null hypothesis.
		\vspace{0.3cm}
		The regression results show that both the number of participating countries and the level of sanctions are associated with individuals’ support for the policy. In general, policies with greater international participation and stronger sanctions tend to increase the likelihood that respondents support the agreement. This suggests that individuals prefer climate agreements that are broadly adopted and credibly enforced.
			\vspace{0.3cm}
	
	\textbf{Conclusion:}
	\textbf{There is strong statistical evidence that policy design influences public support for international environmental agreements. Specifically, increasing participation among countries and implementing sanctions for non-compliance significantly affect the probability that an individual supports the policy. Thus, international cooperation and enforcement mechanisms appear to play an important role in shaping public approval of climate policies}
		
		

	\end{enumerate}
	
	\item
	If any of the explanatory variables are statistically significant in this model, then:
	\begin{enumerate}
		\item
		For the policy in which nearly all countries participate [160 of 192], how does increasing sanctions from 5\% to 15\% change the odds that an individual will support the policy? (Interpretation of a coefficient)
			\lstinputlisting[language=R, firstline=70,lastline=94]{PS02_RS.R} 
			\vspace{10cm}
			\begin{verbatim}
				Call:
				glm(formula = choice ~ countries + sanctions, family = binomial,
				data = climateSupport)
				
				Coefficients:
				Estimate Std. Error z value Pr(>|z|)
				(Intercept) -0.005665   0.021971  -0.258 0.796517
				countries.L  0.458452   0.038101  12.033  < 2e-16 ***
				countries.Q -0.009950   0.038056  -0.261 0.793741
				sanctions.L -0.276332   0.043925  -6.291 3.15e-10 ***
				sanctions.Q -0.181086   0.043963  -4.119 3.80e-05 ***
				sanctions.C  0.150207   0.043992   3.414 0.000639 ***
				---
				Signif. codes:
				0 ‘***’ 0.001 ‘**’ 0.01 ‘*’ 0.05 ‘.’ 0.1 ‘ ’ 1
				
				(Dispersion parameter for binomial family taken to be 1)
				
				Null deviance: 11783  on 8499  degrees of freedom
				Residual deviance: 11568  on 8494  degrees of freedom
				AIC: 11580
				
				Number of Fisher Scoring iterations: 4
			\end{verbatim}
			\textbf{Coefficient:         Odds Ratio = 2, Log-Odds Difference  =
				0.7224531}
\textbf{For policies in which nearly all countries participate (160 of 192), increasing sanctions from 5 percent to 15 percent multiplies the odds that an individual supports the policy by approximately 2. In other words, the odds of supporting the policy are about twice as large when sanctions increase from 5 percent  to 15 percent,  holding the number of participating countries constant.}
			
		
		\item
		For the policy in which very few countries participate [20 of 192], how does increasing sanctions from 5\% to 15\% change the odds that an individual will support the policy? (Interpretation of a coefficient)
			\lstinputlisting[language=R, firstline=98,lastline=111]{PS02_RS.R} 
			\textbf{Coefficient:         Odds Ratio = 2, Log-Odds Difference  =
				0.7224531}
		\item
		What is the estimated probability that an individual will support a policy if there are 80 of 192 countries participating with no sanctions? 
		\lstinputlisting[language=R, firstline=112,lastline=123]{PS02_RS.R} 
		\textbf{Coefficient:         Odds Ratio = 1, Log-Odds Difference  = 0.5159191 }
	
	\end{enumerate}
	\item
	Would the answers to 2a and 2b potentially change if we included an interaction term in this model? Why? 
	\begin{itemize}
		\item Perform a test to see if including an interaction is appropriate.
			\lstinputlisting[language=R, firstline=126,lastline=135]{PS02_RS.R} 
				\vspace{10cm}
			\begin{table}[htbp]
				\centering
				\caption{Logistic Regression Results: Choice $\sim$ Countries $\times$ Sanctions}
				\begin{tabular}{lrrrr}
					\toprule
					& Estimate & Std. Error & z value & Pr($>|z|$) \\
					\midrule
					(Intercept) & -0.0038 & 0.0220 & -0.173 & 0.8626 \\
					
					countries.L & 0.4571 & 0.0381 & 12.005 & $< 2\times10^{-16}$*** \\
					countries.Q & -0.0112 & 0.0382 & -0.293 & 0.7698 \\
					
					sanctions.L & -0.2742 & 0.0440 & -6.239 & $4.41\times10^{-10}$*** \\
					sanctions.Q & -0.1823 & 0.0440 & -4.142 & $3.45\times10^{-5}$*** \\
					sanctions.C & 0.1532 & 0.0441 & 3.477 & 0.000506*** \\
					
					countries.L:sanctions.L & -0.0018 & 0.0767 & -0.023 & 0.9818 \\
					countries.Q:sanctions.L & 0.1338 & 0.0756 & 1.771 & 0.0765. \\
					countries.L:sanctions.Q & -0.0076 & 0.0762 & -0.100 & 0.9203 \\
					countries.Q:sanctions.Q & 0.0934 & 0.0763 & 1.224 & 0.2208 \\
					countries.L:sanctions.C & 0.0952 & 0.0756 & 1.259 & 0.2080 \\
					countries.Q:sanctions.C & 0.0104 & 0.0770 & 0.136 & 0.8921 \\
					\midrule
					\multicolumn{5}{l}{\textit{Model statistics}} \\
					Null deviance & \multicolumn{4}{r}{11783 (df = 8499)} \\
					Residual deviance & \multicolumn{4}{r}{11562 (df = 8488)} \\
					AIC & \multicolumn{4}{r}{11586} \\
					\bottomrule
				\end{tabular}
				
				\vspace{0.5em}
				\begin{minipage}{0.9\linewidth}
					\footnotesize
					\textit{Notes:} Logistic regression with binomial link. 
					Significance levels: *** $p<0.001$, ** $p<0.01$, * $p<0.05$, . $p<0.1$.
				\end{minipage}
			\end{table}
			
			\begin{table}[htbp]
				\centering
				\caption{Analysis of Deviance Table}
				\begin{tabular}{lrrrrr}
					\toprule
					Model & Resid.\ Df & Resid.\ Dev & Df & Deviance & Pr($>\chi^2$) \\
					\midrule
					1: choice $\sim$ countries + sanctions 
					& 8494 & 11568 &  &  &  \\
					
					2: choice $\sim$ countries $\times$ sanctions 
					& 8488 & 11562 & 6 & 6.2928 & 0.3912 \\
					\bottomrule
				\end{tabular}
				
				\vspace{0.5cm}
				\begin{minipage}{0.9\linewidth}
					\footnotesize
					\textit{Notes:} Likelihood ratio test comparing nested logistic regression models.
					Model 2 includes the interaction between countries and sanctions.
				\end{minipage}
			\end{table}
			
			\textbf{Yes, the answers to 2(a) and 2(b) could change if an interaction term were included. In an additive model, the effect of sanctions is assumed to be the same for all levels of country participation, so increasing sanctions from 5 percent to 15 percent changes the odds of support by the same amount regardless of how many countries participate. An interaction model allows the effect of sanctions to vary depending on participation levels.}
			
			\textbf{A likelihood ratio test comparing the additive and interaction models gives a p-value of 0.3912. Since this value is greater than 0.05, we fail to reject the null hypothesis that the interaction terms are zero. Therefore, including an interaction is not appropriate, and the additive model is sufficient. Consequently, the answers to 2(a) and 2(b) remain the same.}
			
	\end{itemize}
	\end{enumerate}


\end{document}
